% =====================================================================
\section{光锥序中间态的取向与 \texorpdfstring{$k^-$}{k-} 极点}\label{sec:dy-orientation}
% =====================================================================

\cref{prop:dy-glauber-pinch}是逐条线算出来的：旁观者的极点在上半平面，活动部分子的极点在下半平面。这种“逐线极点表”在多圈图中很快变得无法管理，而且它掩盖了真正起作用的结构。正确的表述在光锥序微扰论（light-cone ordered perturbation theory，LCOPT）中：对 $J_A$ 按光锥时间 $\tau=x^+$ 排序，软动量的负分量 $k^-$ 是注入的“光锥能量”；它出现在哪些能量分母中，由中间态位于硬顶点之前还是之后决定，与它附着在哪条 Feynman 线上无关。本节给出这一表述；附录中的精确形式与完整证明见\cref{sec:appO-orientation}。

先回顾所需的记号（\cref{def:appO-lcopt}）。固定 $J_A$ 中全部线的正分量与横分量，且每条线 $l^+\neq0$。在壳光锥能量 $E(l)\equiv(\mathbf l_T^2+m_l^2-\ii0)/(2l^+)$ 使 $l^2-m_l^2+\ii0=2l^+\big(l^--E(l)\big)$；对 $l^+>0$，$E(l)$ 位于下半平面。单线的时间表示
\begin{equation}\label{eq:dy-lc-fourier}
\int\frac{\dd l^-}{2\pi}\,\ee^{-\ii l^-\tau}\,\frac{\ii}{2l^+\big(l^--E(l)\big)}=\frac{\theta(l^+)\theta(\tau)+\theta(-l^+)\theta(-\tau)}{2|l^+|}\,\ee^{-\ii E(l)\tau}
\end{equation}
表明 $l^+>0$ 的线沿 $x^+$ 向前传播。把 $J_A$ 的全部负分量积掉，得到对与线的方向相容的顶点排序 $\pi$ 的求和，每个排序给出能量分母之积 $\prod_i\ii/D^\pi_i$，$D^\pi_i=Q^\pi_i-\sum_{l\in S^\pi_i}E(l)$，其中 $S^\pi_i$ 是第 $i$ 个中间态的线集，$Q^\pi_i$ 是到第 $i$ 个顶点为止注入的负分量之和（\cref{lem:appO-lcopt}，\tagR，$R_{\mathrm{PT}}$，在有限的 $\ii\eta_0$ 下成立，并可取 $\eta_0\to0^+$ 的分布极限）。在 DY 的割图中，振幅一侧硬顶点 $V_H$ 之前的中间态称为\emph{初态}，$V_H$ 之后、割线之前的称为\emph{末态}，共轭一侧为镜像（\cref{def:appO-orientation}）。软动量 $k_j$ 的\emph{标准流}是：$k_j$ 在其附着点进入 $J_A$，沿同一侧的路径流到 $V_H$（或 $V_H^\dagger$），经活动线离开。

\begin{lemma}[光锥序中间态的取向与 $k^-$ 极点 \tagR]\label{lem:dy-poles-orientation}
（$R_{\mathrm{PT}}$：在重整化微扰论中对红外正规化的部分子态严格。）本引理只涉及 Feynman 积分。考虑 $J_A$ 在有限 $\eta_0$ 下的光锥序表示（\cref{lem:appO-lcopt}），其活动动量在一个顶点 $V_H$ 处离开振幅部分，软动量取标准流。
\begin{enumerate}[label=(\alph*),leftmargin=2.2em,itemsep=1pt,topsep=2pt]
\item 取向是光锥序中间态的性质，而不是 Feynman 线的性质。振幅一侧的每个初态分母形如 $D=A_D+\sum_{j\in I(D)}k_j^-+\ii0$，每个末态分母形如 $D=A'_D-\sum_{j\in F(D)}k_j^-+\ii0$；这里 $A_D$、$A'_D$ 与全部软动量的负分量无关，$I(D)$ 是在该中间态之前注入的软动量的集合，$F(D)$ 是在其之后、割线之前注入的集合。除此之外，$k_j^-$ 只以多项式形式出现在顶点与接触因子中，不出现在 $1/(2l^+)$、在壳分子与割线的 $\delta$ 函数中。同一条旁观者线可以既属于初态又属于末态。
\item 因此，作为 $k_j^-$ 的函数，初态分母的零点位于 $\Imag k_j^-<0$，末态分母的零点位于 $\Imag k_j^->0$。耦合到软动量 $k$ 的线的 $k^-$ 极点全在同一侧，当且仅当 $k$ 所经过的全部中间态具有同一取向。共轭一侧镜像；对 $J_B$，交换 $+\leftrightarrow-$ 并以 $-k_j$ 为注入动量，$k_j^+$ 有同样的结论。
\end{enumerate}
\uses{lem:appO-lcopt,lem:coll-final-poles,prop:dy-glauber-pinch,def:appO-lcopt,def:appO-orientation,lem:appO-orientation}
\end{lemma}

\begin{proof}
这是\cref{lem:appO-orientation}（\tagR，$R_{\mathrm{PT}}$）的章内表述，完整证明（包括瞬时项与接触顶点）见\cref{sec:appO-orientation}；核心如下。在 $V_H$ 之前的第 $i$ 个中间态，注入的负分量是强子的 $P_A^-$ 加上此前注入的软能量，$D_i=P_A^-+\sum_{j\in I(D_i)}k_j^--\sum_{l\in S_i}E(l)$。在 $V_H$ 之后，振幅部分负分量的总守恒 $P_A^-+K_{\rm tot}-u^-=\mathcal E$（$u$ 为离开 $V_H$ 的活动动量，$\mathcal E=\sum_{\rm cut}E(p_c)$ 为割线上各线在壳能量之和）给出第 $i$ 个末态的注入能量 $\mathcal E-K_{>i}$（$K_{>i}$ 为该态之后、割线之前注入的软能量）。$E(l)$ 只依赖正分量与横分量，割线的 $\delta$ 函数只约束 $\mathcal E$，都不含软负分量；两种情形中 $\Imag D=\sum_{l\in S}\eta_0/(2l^+)>0$。这就是 (a)；(b) 中零点的位置由此直接读出。若 $k$ 只经过初态，$k^-$ 的围道可以整体推入上半平面（这只增加每个初态分母的虚部）；只经过末态的情形就是\cref{lem:coll-final-poles}（\tagR，$R_{\mathrm{PT}}$）；\cref{prop:dy-glauber-pinch}中的 $s_A$ 在一个光锥序项中属于初态（胶子在 $V_H$ 之前附着）、在另一项中属于末态，因而极点分居两侧。
\end{proof}

\cref{fig:dy-pole-orientation}画出了同一个费曼图的两个光锥序项：软胶子在 $V_H$ 之前附着于 $s_A$ 时只进入初态分母，在 $V_H$ 之后附着时只进入末态分母，二者之和使 $k^-$ 围道被钳住。还应注意 (b) 中“不同项之间的钳夹”的意义：Ward 恒等式只在对附着位置求和之后成立，所以全部光锥序项必须共用一条围道；位于向上与向下两种形变之间的极点（例如纵向胶子项在 $k^-=0$ 处的极点）的留数正是 Glauber 相位（\cref{prop:method-glauber-oneloop}，\tagR，$R_{\mathrm{PT}}$）。

%%% BEGIN fig:dy-pole-orientation
\begin{figure}[htbp]
\centering
\begin{tikzpicture}[fd,
  fc hadron/.style={draw=black!70,line width=1.5pt,line cap=butt},
  fc soft/.style={gluon,draw=green!45!black},
  fc state/.style={draw=gray!70!black,line width=0.6pt,dash pattern=on 0.6pt off 1.6pt},
  fc stateB/.style={draw=gray!62,line width=0.6pt,dash pattern=on 0.6pt off 1.6pt},
  fc ax/.style={lcaxis,gray!60!black},
  fc lab/.style={font=\small},
  fc note/.style={font=\footnotesize},
  fc tiny/.style={font=\scriptsize},
  fc pole/.pic={\draw[line width=0.9pt,line cap=round] (-0.075,-0.075)--(0.075,0.075) (-0.075,0.075)--(0.075,-0.075);},
  fc contour/.style={draw=blue!65!black,line width=1.05pt,postaction={decorate,decoration={markings,
      mark=at position #1 with {\arrow{Stealth[length=5pt,width=4.2pt]}}}}},
  fc contour/.default=0.5,
  fc band/.style={fill=red!9,draw=none},
  fc bandA/.style={fill=blue!5,draw=none},
  fc bandB/.style={fill=orange!8,draw=none}]
  % =====================================================================
  % (a) the soft gluon is attached before V_H (tau<0) on both spectators
  % =====================================================================
  \begin{scope}
    % jet subgraphs J_A (top, ordered in x^+) and J_B (bottom, ordered in x^-)
    \fill[fc bandA] (-2.75,0.27) rectangle (2.22,1.95);
    \fill[fc bandB] (-2.75,-1.95) rectangle (2.22,-0.27);
    \node[fc note,text=blue!45!black,anchor=north east] at (2.17,1.5) {$J_A$};
    \node[fc note,text=orange!45!black,anchor=south east] at (2.17,-1.5) {$J_B$};
    % light-cone time axes: x^+ for J_A (top), x^- for J_B (bottom)
    \draw[fc ax] (-3.55,2.75) -- (2.75,2.75) node[right,fc note,text=black] {$x^+$};
    \draw[fc ax] (-3.55,-2.75) -- (2.75,-2.75) node[right,fc note,text=black] {$x^-$};
    \draw[gray!60!black,line width=0.5pt] (0,2.68) -- (0,2.82) (0,-2.68) -- (0,-2.82);
    \node[fc tiny,above] at (0,2.78) {$0$};  \node[fc tiny,below] at (0,-2.78) {$0$};
    \node[fc tiny,above,text=gray!40!black] at (-1.75,2.78) {初态};
    \node[fc tiny,above,text=gray!40!black] at (1.2,2.78) {末态};
    \node[fc tiny,below,text=gray!40!black] at (-1.75,-2.78) {初态};
    \node[fc tiny,below,text=gray!40!black] at (1.2,-2.78) {末态};
    % intermediate states (vertical dotted lines)
    \draw[fc state] (-2.0,1.05) -- (-2.0,1.85);
    \draw[fc state] (-0.7,0.22) -- (-0.7,1.85);
    \draw[fc stateB] (-2.0,-1.05) -- (-2.0,-1.85);
    \draw[fc stateB] (-0.7,-0.22) -- (-0.7,-1.85);
    \node[fc tiny,text=gray!40!black] at (-2.0,2.08) {$D_1$};
    \node[fc tiny] at (-0.15,2.27) {$\dfrac{1}{D_2}=\dfrac{1}{k^-+A+\ii0}$};
    % soft gluon k from s_B to s_A (drawn first: it passes behind a and b)
    \draw[fc soft] (-1.5,-1.6) -- (-1.5,1.6);
    \path[mom'={$k$},mom pos=0.5,mom sep=0.27cm] (-1.5,-1.6) -- (-1.5,1.6);
    % hadrons, splitting vertices, spectators, active partons
    \draw[fc hadron] (-3.6,1.6) -- (-2.5,1.6);
    \draw[fc hadron] (-3.6,-1.6) -- (-2.5,-1.6);
    \draw[plain] (-2.5,1.6) -- (2.75,1.6);
    \draw[plain] (-2.5,-1.6) -- (2.75,-1.6);
    % a and b pass in front of the soft gluon (gap in the background colour of the band)
    \draw[blue!5,line width=3.8pt] (-1.71,1.095) -- (-1.29,0.825);
    \draw[orange!8,line width=3.8pt] (-1.71,-1.095) -- (-1.29,-0.825);
    \draw[quark=0.62] (-2.5,1.6) -- (-0.17,0.11);
    \draw[quark=0.38] (-0.17,-0.11) -- (-2.5,-1.6);
    \draw[photon] (0.2,0) -- (0.9,0);
    \node[hard blob,circle,minimum size=0.4cm,inner sep=0pt] (H) at (0,0) {};
    \foreach \p in {(-2.5,1.6),(-2.5,-1.6),(-1.5,1.6),(-1.5,-1.6)} \node[vertex] at \p {};
    \draw[cut] (2.3,-2.05) -- (2.3,2.05);
    % labels
    \node[fc lab,above] at (-3.25,1.62) {$P_A$};
    \node[fc lab,below] at (-3.25,-1.62) {$P_B$};
    \node[fc lab,right] at (2.75,1.6) {$s_A$};
    \node[fc lab,right] at (2.75,-1.6) {$s_B$};
    \node[fc lab] at (-2.25,0.98) {$a$};
    \node[fc lab] at (-2.25,-0.98) {$b$};
    \node[fc note,anchor=north west] at (0.12,-0.12) {$V_H$};
    \node[fc note,above] at (0.62,0.07) {$\gamma^*$};
    \node[fc lab] at (-0.4,-3.6) {(a)\ $k$在$V_H$之前附着：只进入初态};
  \end{scope}
  % =====================================================================
  % (b) the same Feynman graph, soft gluon attached after V_H (tau>0)
  % =====================================================================
  \begin{scope}[xshift=7.45cm]
    \fill[fc bandA] (-2.75,0.27) rectangle (2.22,1.95);
    \fill[fc bandB] (-2.75,-1.95) rectangle (2.22,-0.27);
    \node[fc note,text=blue!45!black,anchor=north east] at (2.17,1.5) {$J_A$};
    \node[fc note,text=orange!45!black,anchor=south east] at (2.17,-1.5) {$J_B$};
    \draw[fc ax] (-3.55,2.75) -- (2.75,2.75) node[right,fc note,text=black] {$x^+$};
    \draw[fc ax] (-3.55,-2.75) -- (2.75,-2.75) node[right,fc note,text=black] {$x^-$};
    \draw[gray!60!black,line width=0.5pt] (0,2.68) -- (0,2.82) (0,-2.68) -- (0,-2.82);
    \node[fc tiny,above] at (0,2.78) {$0$};  \node[fc tiny,below] at (0,-2.78) {$0$};
    \node[fc tiny,above,text=gray!40!black] at (-1.75,2.78) {初态};
    \node[fc tiny,above,text=gray!40!black] at (1.2,2.78) {末态};
    \node[fc tiny,below,text=gray!40!black] at (-1.75,-2.78) {初态};
    \node[fc tiny,below,text=gray!40!black] at (1.2,-2.78) {末态};
    \draw[fc state] (-1.25,0.58) -- (-1.25,1.85);
    \draw[fc state] (0.62,1.35) -- (0.62,1.85);
    \draw[fc stateB] (-1.25,-0.58) -- (-1.25,-1.85);
    \draw[fc stateB] (0.62,-1.35) -- (0.62,-1.85);
    \node[fc tiny,text=gray!40!black] at (-1.25,2.08) {$D_1$};
    \node[fc tiny] at (0.85,2.27) {$\dfrac{1}{D_2}=\dfrac{-1}{k^--A'-\ii0}$};
    \draw[fc soft] (1.25,-1.6) -- (1.25,1.6);
    \path[mom'={$k$},mom pos=0.5,mom sep=0.27cm] (1.25,-1.6) -- (1.25,1.6);
    \draw[fc hadron] (-3.6,1.6) -- (-2.5,1.6);
    \draw[fc hadron] (-3.6,-1.6) -- (-2.5,-1.6);
    \draw[plain] (-2.5,1.6) -- (2.75,1.6);
    \draw[plain] (-2.5,-1.6) -- (2.75,-1.6);
    \draw[quark=0.7] (-2.5,1.6) -- (-0.17,0.11);
    \draw[quark=0.3] (-0.17,-0.11) -- (-2.5,-1.6);
    \draw[photon] (0.2,0) -- (0.9,0);
    \node[hard blob,circle,minimum size=0.4cm,inner sep=0pt] (H) at (0,0) {};
    \foreach \p in {(-2.5,1.6),(-2.5,-1.6),(1.25,1.6),(1.25,-1.6)} \node[vertex] at \p {};
    \draw[cut] (2.3,-2.05) -- (2.3,2.05);
    \node[fc lab,above] at (-3.25,1.62) {$P_A$};
    \node[fc lab,below] at (-3.25,-1.62) {$P_B$};
    \node[fc lab,right] at (2.75,1.6) {$s_A$};
    \node[fc lab,right] at (2.75,-1.6) {$s_B$};
    \node[fc lab] at (-2.25,0.98) {$a$};
    \node[fc lab] at (-2.25,-0.98) {$b$};
    \node[fc note,anchor=north west] at (0.12,-0.12) {$V_H$};
    \node[fc note,above] at (0.55,0.07) {$\gamma^*$};
    \node[fc lab] at (-0.4,-3.6) {(b)\ $k$在$V_H$之后附着：只进入末态};
  \end{scope}
  % =====================================================================
  % (c), (d): complex k^- plane
  % =====================================================================
  \begin{scope}[shift={(1.2,-5.45)}]
    \fill[fc band] (-0.42,-0.85) rectangle (0.42,0.85);
    \draw[fc ax] (-1.75,0) -- (1.8,0) node[below,fc tiny,text=black,xshift=-3pt] {$\Real k^-$};
    \draw[fc ax] (0,-0.9) -- (0,1.0) node[left,fc tiny,text=black] {$\Imag k^-$};
    \draw[fc contour=0.85] (-1.65,0) -- (1.65,0);
    \pic at (-0.16,-0.26) {fc pole};  \pic at (0.16,0.26) {fc pole};
    \node[fc tiny,anchor=east] at (-0.24,-0.34) {(a)};
    \node[fc tiny,anchor=west] at (0.24,0.34) {(b)};
    \node[fc note,align=center] at (0,-1.52) {(c)\ (a)+(b)：极点分居两侧，\\$k^-$围道被钳住};
  \end{scope}
  \begin{scope}[shift={(5.7,-5.45)}]
    \fill[fc band] (-0.42,-0.85) rectangle (0.42,0.85);
    \draw[fc ax] (-1.75,0) -- (1.8,0) node[below,fc tiny,text=black,xshift=-3pt] {$\Real k^-$};
    \draw[fc ax] (0,-0.9) -- (0,1.0) node[left,fc tiny,text=black] {$\Imag k^-$};
    \draw[fc contour=0.5] (-1.65,0) -- (-1.2,0) .. controls (-0.95,0) and (-0.85,0.62) .. (0,0.62)
        .. controls (0.85,0.62) and (0.95,0) .. (1.2,0) -- (1.65,0);
    \pic at (-0.2,-0.22) {fc pole};  \pic at (0.12,-0.3) {fc pole};
    \node[fc note,align=center] at (0,-1.52) {(d)\ 只经过初态：极点同侧，\\$k^-$围道可移入上半平面};
  \end{scope}
\end{tikzpicture}
\caption{光锥序中间态的取向与$k^-$极点（\cref{lem:dy-poles-orientation}）。Drell--Yan 割图的振幅一侧（割线右侧的共轭振幅为镜像，未画出）按混合光锥序排列：$J_A$（上，浅蓝）按$x^+$、$J_B$（下，浅橙）按$x^-$排序，硬顶点$V_H$在两者中都位于$\tau=0$；
竖直点线标出中间态，$a$、$b$为活动部分子，$s_A$、$s_B$为旁观者，软胶子（绿色）动量$k$由$J_B$流入$J_A$（标准流）。
(a)、(b) 是同一费曼图的两个光锥序项：$k$在$V_H$之前附着于$s_A$时，$k^-$只出现在初态分母$D_2=k^-+A+\ii0$中，零点在下半平面；在$V_H$之后附着时，$k^-$只出现在末态分母$D_2=-(k^--A'-\ii0)$中，零点在上半平面
（$A$、$A'$为与$k^-$无关的实数；$J_B$中对$k^+$有同样的结论，$+\leftrightarrow-$）。取向因而是中间态而不是线的性质：(c)~$k$流经两种取向的中间态时，极点分居实轴两侧，$k^-$围道被钳在 Glauber 区（阴影）内，旁观者之间的交换在$J_A$、$J_B$中同时被钳住，即 Glauber 钳夹；
(d)~$k$只流经初态（例如只附着于$a$）时极点同侧，$k^-$围道可移入上半平面（离开 Glauber 区还须同时形变$k^+$，见\cref{fig:method-glauber-direction}）。}
\label{fig:dy-pole-orientation}
\end{figure}
%%% END fig:dy-pole-orientation

\begin{remark}[每个喷注一个光锥时间]\label{rem:dy-one-time-per-jet}
上面的构造对 $J_A$ 用光锥时间 $x^+$（能量为负分量），对 $J_B$ 用 $x^-$（能量为正分量），硬顶点在两者中都位于 $\tau=0$：这是 Aybat 与 Sterman 引入的“混合”光锥序\cite{AybatSterman2009}。选择的理由在于：$A$ 共线线的大分量是正分量，以 $x^+$ 为时间时，Glauber 与软动量只以小的能量 $k^-$ 进入其能量分母；在 Glauber 区内连接线的传播子 $\simeq-1/\mathbf k_T^2$ 与 $k^\pm$ 无关，因而在坐标空间中正比于 $\delta(x^+-y^+)\delta(x^--y^-)$，而 $J_A$ 不感知 $k^+$（附着点在 $x^-\simeq0$）、$J_B$ 不感知 $k^-$（附着点在 $y^+\simeq0$）——Glauber 交换发生在硬顶点的光锥时刻，即两个被收缩的强子彼此穿过的瞬间。后文所说的“测量与光锥演化相容”，精确的意思是：测量算符在割线两侧的末态之间作为恒等算符（或与演化对易的算符）出现；其一般判据见\cref{prop:break-measurement}（\tagR，$R_{\mathrm{cond}}$）。
\end{remark}
